\documentclass[11pt,a4paper,oneside,fontset=windows]{ctexbook}
\usepackage{amsthm,amssymb} % math packages
\usepackage{amsmath}
\usepackage{mathrsfs}
\usepackage{bm} % bold math
\usepackage[hidelinks]{hyperref} % hide ref in table of contents
\usepackage{pifont} % \ding{73} makes a beautiful star
\usepackage{tikz}
\usepackage{fontspec}
\usepackage{color}
\usepackage{geometry}
\usepackage{stmaryrd}
\usepackage{subfiles}
\usepackage{pgfplots}
\usepackage{pgfornament}
\usepackage{enumerate}
\usepackage{mdwlist}
\usepackage{booktabs}
\usepackage{ulem}
\usepackage[level]{datetime}


\allowdisplaybreaks[4]

\setcounter{tocdepth}{1}
\definecolor{navy}{rgb}{0,0,0.5}
% \tag{} numbers the equation
\def\d#1{\text{d}#1}
\def\D#1{\text{d}(#1)}
\def\p#1{\partial#1}
\def\ds#1#2{\frac{\text{d}#1}{\text{d}#2}}
\def\pds#1#2{\frac{\partial#1}{\partial#2}}
\def\jishu#1#2{\sum_{#1=1}^{+\infty}#2}
\def\linsp#1{\mathbb{R}^{#1}}
\def\ntoinf#1{\lim\limits_{#1\rightarrow\infty}}
\def\tinysec#1{\noindent\textbf{#1}\indent}
\def\concept#1{\uline{\textbf{#1}}}
\geometry{a4paper, scale=0.75}

\newenvironment{solution}{\noindent{\kaishu 解}.}{  }

\title{\Huge{\heiti 数学支配着宇宙}\\Math Dominates the Universe}
\author{Jingze Xing}

\pgfplotsset{compat=1.16} % advised setting
\begin{document}
%\ThisCenterWallPaper{}{back.jpg}
\thispagestyle{empty}
\begin{tikzpicture}[color=navy,overlay,remember picture,
    every node/.style={inner sep=5pt}]
    \node[anchor=north west] at (current page.north west)
    {\pgfornament[width=3cm]{63}};
    \node[anchor=north east] at (current page.north east)
    {\pgfornament[width=3cm,symmetry=v]{63}};
    \node[anchor=south west] at (current page.south west)
    {\pgfornament[width=3cm,symmetry=h]{63}};
    \node[anchor=south east] at (current page.south east)
    {\pgfornament[width=3cm,symmetry=c]{63}};
\end{tikzpicture}
\begin{figure}
    \centering
    \begin{minipage}[t][10cm][b]{\textwidth}
        \centering
        \begin{tikzpicture}[color=black,remember picture]
            \node[minimum width=12cm, minimum height=3.6cm](vecbox){};
            \node[anchor=north] at (vecbox.north)
            {\pgfornament[height=0.27cm]{89}};
            \node[anchor=south] at (vecbox.south)
            {\pgfornament[height=0.27cm]{89}};
            
            \node[inner sep=10pt,align=center] (text) at (vecbox.center){\Huge\textbf{Math Dominates the Universe}\\[1em]\Huge{\heiti 数学支配着宇宙}};
        \end{tikzpicture}
    \end{minipage}
    \centering
    \begin{minipage}[t][3.5cm][b]{\textwidth}
        \centering
        \begin{tikzpicture}[color=black,remember picture]
            \node[minimum width=8.6cm, minimum height=2cm](vecbox){};
            \node[inner sep=0pt,align=center,anchor=north] (text) at ([yshift={-7pt}]vecbox.north){\LARGE\textbf{Jingze Xing}};
            \node[inner sep=0pt,align=center,anchor=south] (text) at (vecbox.south){\Large{\today}};
            \node[inner sep=0pt,anchor=north west] at (vecbox.north west)
            {\pgfornament[height=1cm]{113}};
            \node[inner sep=0pt,anchor=north east] at (vecbox.north east)
            {\pgfornament[height=1cm,symmetry=v]{113}};
        \end{tikzpicture}
    \end{minipage}
\end{figure}
\tableofcontents
\newpage

\subfile{1-calculus1.tex}
\subfile{2-calculus2.tex}
\subfile{7-linear-algebra.tex}
% \subfile{3-discrete.tex}
\subfile{4-complex-functions.tex}
\subfile{5-probability.tex}
\subfile{6-stochastic-process.tex}

\chapter{线性代数}
\section{基本事实}
\begin{itemize}
  \item \( rank(A)+rank(B)<rank(AB)-n,A\in \mathbb{R}^{m\times n} \)
  \item  \( a_1+a_2,a_2+a_3,\cdots,a_n+a_1 \text{可逆}\implies n\text{为奇数} \)
  \item \( Q^TQ=I_n\iff Q\text{为保距变换}\iff \text{为保内积变换}\)
  \item \textit{Housholder}变换(关于\textbf{面}反射)：\( I_n-2\bm{vv}^T,||\bm{v}||=1 \)
  \item \( \mathcal{R}(A^TA)=\mathcal{R}(A^T)\)
  \item \( P_{\mathcal{M}}=QQ^T,\text{span}(Q)=\mathcal{M} \)
  \item \( P_{A}=A(A^TA)^{-1}A^T,A \)列线性无关
  \item 实对称矩阵\( S=Q\Lambda Q^T \)
\end{itemize}
\section{例题}
\subsubsection{练习1.6.17-(5)}
若 $A^2-I_n=O$ ，证明：
\begin{equation*}
  \forall \bm{v}\in \mathbb{R}^n, \exists \bm{x}+\bm{y}=\bm{v},\ s.t.\ (A-I_n)\bm{x}=(A+I_n)\bm{y}=\bm{0}
\end{equation*}
\begin{proof}
  \begin{align*}
                  & (A+I_n)(A-I_n)=0                                           \\
    \text{构造}\  & \bm{x}=\frac{A+I_n}{2}\bm{v},\bm{y}=\frac{-A+I_n}{2}\bm{v}
  \end{align*}
\end{proof}
\subsubsection{练习 1.7.10}
证明：$A\in\mathbb{R}^{m*n},B\in\mathbb{R}^{n*m}, I_m+AB\ \text{可逆} \iff I_n+BA\ \text{可逆}$.
\begin{proof}
  \begin{align*}
                  & \text{构造矩阵} \begin{bmatrix} I_m  &A \\ B &I_n \end{bmatrix}                       \\
                  & \begin{bmatrix} I_m  &O \\-B  &I_n\end{bmatrix}
    \begin{bmatrix} I_m  &A \\  B  &I_n\end{bmatrix}=
    \begin{bmatrix} I_m  &A \\  O  &I_n-BA\end{bmatrix}                                                       \\
                  & \begin{bmatrix} I_m  &-A \\ O  &I_n\end{bmatrix}
    \begin{bmatrix} I_m  &A \\ B  &I_n\end{bmatrix}=
    \begin{bmatrix} I_m-AB  &O \\ B  &I_n\end{bmatrix}                                                       \\
                  & I_n-BA\text{可逆}\iff
    \begin{bmatrix} I_m &A \\ B &I_n\end{bmatrix}\text{可逆} \iff I_m-AB \text{可逆}                   \\
    \text{另：\ } & (I_m+AB)A=A+ABA=A(I_n+BA)                                       \\
    let\          & (I_n+BA)^{-1}=C                                                 \\
    then\         & (I_m+AB)ACB=AB                                                  \\
                  & (I_m+AB)-(I_m+AB)ACB=(I_m+AB)-AB                                \\
                  & (I_m+AB)(I_m-ACB)=I_m\implies (I_m+AB)^{-1}=I_m-A(I_n+BA)^{-1}B
  \end{align*}
\end{proof}
\subsection{正交矩阵}
\subsection{特征值与特征向量}
\subsubsection{练习5.2.23}
给定$m$阶方阵$A_1$，$n$阶方阵$A_2$和$m\times n$矩阵$B$，证明若$A_1$和$A_2$没有相同的特征值，则关于$m\times n$矩阵$X$的矩阵方程$A_1X-XA_2=B$有唯一解.\\
\begin{proof}
  因为$X$与$B$规格相同，所以只要证明该变换是单射即可证明是双射.\\
  即证：$A_1X=XA_2$仅有零解
  \begin{align*}
    A_1^2X=(A_1X) & A_2=XA_2^2                      \\
    A_1^3X=(A_1X) & A_2^2=XA_2^3                    \\
                  & \vdots                          \\
    A_1^nX=(A_1X) & A_2^{n-1}=XA_2^n\tag{\ding{72}}
  \end{align*}
  考虑上述等式的线性组合，则
  $$f(A_1)X=Xf(A_2)$$
  恒成立.\\
  不妨令$f=p_A$，则$f(A_1)=O$.令$\lambda_i$为$A_1$的特征值，则$\det(\lambda_iI_n-A_2)\neq 0$，所以
  \begin{align*}
    \det(f(A_2))=\prod\limits_{i=1}^{n}{(\lambda_i I_n-A_2)}\neq 0 \tag{\ding{72}}
  \end{align*}
  故$f(A_2)$可逆，因此$$X=f(A_1)X\left[f(A_2)\right]^{-1}=O$$
\end{proof}
\subsection{实对称矩阵}
对实对称矩阵，$A$正定\\$\iff$$A$的特征值都是正数\\$\iff$$A=TT^T$,$T$可逆\\$\iff$$A=LDL^T$\\$\iff$$A$的顺序主子式都是正数
\newpage

\newpage
\begin{itemize}
  \item gbk (gbk不能解码的特殊字符，用更高级的GB18030)
\end{itemize}

\end{document}